\documentclass[11pt, a4paper]{amsart}
\usepackage{amsmath, amssymb, amsthm, amsfonts}
\usepackage{mathtools}
\usepackage[margin=1in]{geometry}
\usepackage{hyperref}
\usepackage{tikz-cd}

\newtheorem{theorem}{Theorem}[section]
\newtheorem{lemma}[theorem]{Lemma}
\newtheorem{proposition}[theorem]{Proposition}
\newtheorem{corollary}[theorem]{Corollary}
\newtheorem{fact}[theorem]{Standard Fact}
\theoremstyle{definition}
\newtheorem{definition}[theorem]{Definition}
\newtheorem{remark}[theorem]{Remark}
\newtheorem{notation}[theorem]{Notation}
\newtheorem{example}[theorem]{Example}

\newcommand{\R}{\mathbb{R}}
\newcommand{\Q}{\mathbb{Q}}
\newcommand{\C}{\mathbb{C}}
\newcommand{\Z}{\mathbb{Z}}
\newcommand{\N}{\mathbb{N}}
\newcommand{\HH}{\mathbb{H}}
\newcommand{\OO}{\mathbb{O}}
\newcommand{\id}{\operatorname{id}}
\newcommand{\im}{\operatorname{im}}
\newcommand{\tr}{\operatorname{tr}}
\newcommand{\spec}{\operatorname{spec}}
\newcommand{\op}{\operatorname{op}}
\newcommand{\Sym}{\operatorname{Sym}}
\newcommand{\Met}{\mathrm{met}}
\newcommand{\PF}{\mathrm{PF}}
\newcommand{\src}{\mathrm{src}}
\newcommand{\Pimet}{\Pi_{\Met}}
\renewcommand{\Im}{\mathrm{Im}}

\title{Cascade--Classical Correspondence V:\\
  $D_4$ Metric Projection from Discrete Moment Tensors}
\author{}
\date{}

\begin{document}

\maketitle

\begin{abstract}
This paper puts a metric-side observable on the refinement-spaces
substrate of paper P3 \cite{RefinementSpaces} using the explicit
$D_4$ coordinates of paper P2 \cite{AlgebraicSubstrate} and the
finite-level closure-dynamics state space of paper P4
\cite{ClosureDynamics}. At each refinement level $n \geq 0$ and on
the $D_4$ branch, we define the \emph{discrete moment tensor}
\[
  M_n[s](v) \;:=\; \sum_{w \in V_n^{D_4},\ w \neq v} s(w)\,
    K_n(v, w), \qquad
  K_n(v, w) \;:=\; \frac{(x_w^{(n)} - x_v^{(n)}) \otimes
      (x_w^{(n)} - x_v^{(n)})}{\|x_w^{(n)} - x_v^{(n)}\|^2},
\]
and the \emph{metric projection}
$\Pimet^{(n)} := P_n^\PF \circ M_n \circ S_n^\lambda \colon X_n^{D_4} \to
\PF_n$ from the P4 state space $X_n^{D_4}$ to the piecewise-flat
target $\PF_n$. We prove six theorems at finite level:

\begin{enumerate}
\item[(T1)] $M_n$ is a well-defined bounded linear map from
  $\ell^2(V_n^{D_4})$ to
  $\ell^2(V_n^{D_4}; \Sym^2(\R^4))$ with each
  $M_n[s](v)$ symmetric (Theorem~\ref{thm:mn-welldef}).
\item[(T2)] On the subspace
  $\mathcal S_{n+1}^{\mathrm{old}}$ of level-$(n+1)$ source
  functions supported on old vertices $V_n \subset V_{n+1}$, the
  raw moment operator satisfies the operator identity
  $q_{n+1, n} \circ M_{n+1} = M_n \circ p_{n+1, n}^0$, where
  $p_{n+1, n}^0$ is the P3 vertex bonding and $q_{n+1, n}$ is the
  vertex restriction on tensor fields defined here
  (Theorem~\ref{thm:mn-refinement}). The identity does not extend
  to all of $\mathcal S_{n+1}$ because new-vertex source data at
  level $n+1$ produces additional kernel contributions.
\item[(T3)] $M_n$ and $\Pimet^{(n)}$ are $W(D_4)$-equivariant
  under the explicit orthogonal action on coordinates coming
  from P2 and the branchwise Coxeter action of P3
  (Theorem~\ref{thm:mn-equivariance}).
\item[(T4)] $\Pimet^{(n)}$ is bounded with
  $\|\Pimet^{(n)}\|_\op \leq \|P_n^\PF\|_\op \cdot \|M_n\|_\op
  \cdot \|S_n^\lambda\|_\op$ (Theorem~\ref{thm:pimet-continuity}).
\item[(T5)] $\Pimet^{(n)}$ lands in the piecewise-flat target
  $\PF_n$, defined explicitly as the finite-dimensional Hilbert
  space of cell-labelled symmetric tensor fields on the
  level-$n$ vertex set of the $D_4$ refinement (one constant
  symmetric tensor per vertex, no cross-cell compatibility
  imposed) with the cellwise Frobenius inner product
  (Theorem~\ref{thm:pimet-landing}).
\item[(T6)] For a point source $s = \delta_{w_\ast}$ and any probe
  $v \neq w_\ast$, $M_n[\delta_{w_\ast}](v) = \hat r_{v, w_\ast}^{(n)}
  \otimes \hat r_{v, w_\ast}^{(n)}$ is the unit rank-one projector
  onto the displacement direction, with trace $1$ and spectrum
  $\{1, 0, 0, 0\}$ (Theorem~\ref{thm:point-source}).
\end{enumerate}

The paper is finite-level only. We prove no continuum metric,
Riemannian limit, Lorentzian structure, Einstein equation,
stress-energy, perfect-fluid, hydrodynamic, Navier--Stokes, or
general-relativistic statement. Convergence under refinement,
Gromov--Hausdorff limits, and spectral stability of graph
Laplacians are the target of paper P6 and are out of scope here.
\end{abstract}

\tableofcontents

%=====================================================================
\section{Introduction and Source Audit}
\label{sec:intro}
%=====================================================================

This is the first metric-side paper in the cascade infrastructure
programme. Its load-bearing job is narrow and entirely
finite-dimensional: starting from the finite-level refinement
graphs $G_n^{D_4}$ of P3~\cite{RefinementSpaces}, the explicit
$D_4$ coordinate model of P2~\cite{AlgebraicSubstrate}, and the
P4 state space $X_n^{D_4}$ of \cite{ClosureDynamics}, we build a
bounded linear operator $\Pimet^{(n)} \colon X_n^{D_4} \to \PF_n$
with the four structural properties expected of a metric-side
observable: refinement compatibility, Coxeter equivariance,
boundedness, and a clean rank-one point-source limit. The paper
is designed to be the theorem-grade replacement for the
exploratory metric-projection discussion in
\texttt{papers/cascade-derivation/cascade-gr.md} \S C4 and in
\texttt{papers/cascade-correspondence-foundations/foundations.tex};
those sources are cited only as historical motivation in this
introduction and never as theorem-grade proof support.

\subsection{What this paper proves}

At each refinement level $n \geq 0$ and on the $D_4$ branch, we
define:
\begin{enumerate}
\item the canonical Euclidean coordinates $x_v^{(n)} \in \R^4$ of
  the refinement vertices $V_n^{D_4}$, inherited from P2
  (\S\ref{sec:coords});
\item a bounded scalar source extraction map $S_n^\lambda \colon
  X_n^{D_4} \to \ell^2(V_n^{D_4})$ from the P4 state space
  (\S\ref{sec:source});
\item the normalized rank-one displacement kernel $K_n(v, w)$
  for $v \neq w$ (\S\ref{sec:moment});
\item the raw discrete moment tensor operator
  $M_n \colon \ell^2(V_n^{D_4}) \to \ell^2(V_n^{D_4}; \Sym^2(\R^4))$
  (\S\ref{sec:moment});
\item the piecewise-flat target $\PF_n$ (a finite-dimensional
  real Hilbert space of cell-labelled symmetric tensors on $V_n$
  with the cellwise Frobenius inner product, no cross-cell
  compatibility imposed), and the piecewise-flat projection
  $P_n^\PF = \id$ (\S\ref{sec:pf-target});
\item the metric projection $\Pimet^{(n)} := P_n^\PF \circ M_n
  \circ S_n^\lambda$ (\S\ref{sec:pf-target}).
\end{enumerate}
We prove the six theorems listed in the abstract, all in finite
dimension and all as explicit bounded-linear-map / operator-identity
statements.

\subsection{What this paper does not prove}

P5 is deliberately narrow. We record the exclusions explicitly:
\begin{enumerate}
\item No continuum-limit theorem. We do not claim
  $\lim_{n \to \infty} \Pimet^{(n)}$ converges, in any sense, to
  a continuum metric, metric-measure space, or Riemannian
  manifold. All convergence claims belong to P6.
\item No Gromov--Hausdorff, distortion, $\varepsilon$-net, or
  spectral-stability statement for the refinement family.
\item No Lorentzian, Einstein, Ricci, stress-energy, perfect-fluid,
  conservation-law, or graviton claim.
\item No identification of the finite-level metric-side output
  with any classical metric on $S^3$, $\R^4$, $\R^{1,3}$, or any
  other manifold.
\item No hydrodynamic projection, macroscopic PDE, Euler-equation,
  or Navier--Stokes statement. Those belong to P7.
\item No quotient map $V_n^{H_4} \to V_n^{D_4}$; the paper works
  entirely on the $D_4$ branch.
\item No physical or phenomenological interpretation of the
  tensor fields produced by $\Pimet^{(n)}$.
\end{enumerate}

\subsection{Source audit}

Theorem inputs to this paper are:
\begin{description}
\item[\textbf{Internal cascade infrastructure.}]
P2~\cite{AlgebraicSubstrate} for the explicit $24$-cell vertex set
$V_{24} \subset \R^4$, its Coxeter symmetry $W(D_4)$ acting
orthogonally on $\R^4$, and the identification of $V_{24}$ with
the rescaled $D_4$ root polytope. P3~\cite{RefinementSpaces} for
the refinement graphs $G_n^{D_4}$, the vertex inclusion
$V_n^{D_4} \subset V_{n+1}^{D_4}$, the vertex bonding
$p_{n+1, n}^0$ (restriction), and the equivariance of refinement
under $W(D_4)$. P4~\cite{ClosureDynamics} for the finite-level
state space $X_n^{D_4} = X_n^{0, D_4} \oplus X_n^{1, D_4}$, its
Hilbert structure, and the bounded vertex-projection
$X_n^{D_4} \to X_n^{0, D_4}$. All three papers have been
independently reviewed at Clay-bar rigour and accepted.
\item[\textbf{Classical mathematics.}]
Finite-dimensional linear algebra of symmetric bilinear forms on
$\R^4$ and rank-one tensor products \cite{HornJohnson}; elementary
$\ell^2$ theory on finite sets \cite{ReedSimonI}; standard
properties of orthogonal actions of finite Coxeter groups on
Euclidean space \cite{Humphreys}.
\end{description}
We do not cite as theorem-grade input any other cascade-internal
note. The bibliography below contains only the sources actually
used in proofs.

\subsection{Organization}

Section~\ref{sec:coords} records the level-$n$ vertex coordinates
and the $W(D_4)$ action on $\R^4$. Section~\ref{sec:source}
defines the source space $\mathcal S_n$ and the bounded extraction
map $S_n^\lambda$. Section~\ref{sec:moment} defines $K_n$ and $M_n$ and
proves well-definedness (Theorem~\ref{thm:mn-welldef}).
Section~\ref{sec:pf-target} defines the piecewise-flat target
$\PF_n$, its metric, the projection $P_n^\PF$, and the metric
projection $\Pimet^{(n)}$. Section~\ref{sec:refinement} proves
refinement stability (Theorem~\ref{thm:mn-refinement}).
Section~\ref{sec:equivariance} proves $W(D_4)$-equivariance
(Theorem~\ref{thm:mn-equivariance}).
Section~\ref{sec:continuity} proves continuity with an explicit
operator-norm bound (Theorem~\ref{thm:pimet-continuity}) and the
piecewise-flat landing theorem
(Theorem~\ref{thm:pimet-landing}). Section~\ref{sec:point-source}
proves the point-source rank-one calculation
(Theorem~\ref{thm:point-source}). Section~\ref{sec:audit} is a
self-audit against the Clay-bar success criteria of the work
order. Section~\ref{sec:handoff} gives the citation contract for
downstream papers. Section~\ref{sec:notation} is the notation
audit.

%=====================================================================
\section{The $D_4$ Branch: Coordinates and Symmetry}
\label{sec:coords}
%=====================================================================

Throughout this paper $\bullet = D_4$ unless stated otherwise; we
drop the decoration when no confusion arises.

\begin{notation}[Level-$n$ $D_4$ vertices and coordinates]
\label{not:vertices-coords}
For each $n \geq 0$, let $V_n := V_n^{D_4} = V(G_n^{D_4})$ be the
level-$n$ vertex set of the $D_4$ refinement graph of
\cite[Def.~2.1, Def.~2.2]{RefinementSpaces}. For each
$v \in V_n$, let $x_v^{(n)} \in \R^4$ denote the canonical
Euclidean coordinate of $v$ obtained from the P2
$24$-cell realization \cite[\S 5]{AlgebraicSubstrate} together
with the Schl\"{a}fli-refinement midpoint / centroid rule
\cite[Def.~2.2]{RefinementSpaces}. Under vertex inclusion
\cite[Rmk.~2.3]{RefinementSpaces}, these coordinates are
compatible: $x_v^{(n+1)} = x_v^{(n)}$ for $v \in V_n \subset
V_{n+1}$.
\end{notation}

\begin{hypothesis}[Coordinate-injectivity of the refinement]
\label{fact:distinct}
The coordinate map $V_n \to \R^4$, $v \mapsto x_v^{(n)}$, is
injective for every $n \geq 0$.
\end{hypothesis}

\begin{remark}[Status of Hypothesis~\ref{fact:distinct}]
\label{rmk:distinct-status}
At level $0$, $V_0 = V_{24}$ is the $24$-vertex set of the
$24$-cell in $\R^4$, whose vertices are pairwise distinct
\cite[\S 5]{AlgebraicSubstrate}, so injectivity holds at level
$0$. At higher levels, new vertices are added as edge-midpoints
and face-centroids of level-$n$ objects
\cite[Def.~2.2]{RefinementSpaces}. A level-$(n+1)$ midpoint lies
in the relative interior of the spanning $1$-face and a
level-$(n+1)$ centroid lies in the relative interior of the
spanning $2$-face, so each new coordinate lies strictly inside a
cell of level $n$ and is distinct from every level-$n$ vertex.
However, proving that \emph{distinct} new midpoints/centroids
produce distinct coordinate points, and that midpoint-coordinates
do not coincide with centroid-coordinates of a different $2$-face,
requires a finer Schl\"afli-subdivision-geometry argument beyond
``relative-interior'' which this paper does not supply. We
therefore take Hypothesis~\ref{fact:distinct} as a structural
hypothesis of the refinement scheme; it is satisfied for the
standard Schl\"afli refinement of the $24$-cell by direct
coordinate check at each level, but the full geometric argument
establishing it for all $n$ is deferred. All subsequent results
of this paper are stated conditional on
Hypothesis~\ref{fact:distinct}.
\end{remark}

\begin{corollary}[Nonvanishing denominator]
\label{cor:nonvanishing}
Under Hypothesis~\ref{fact:distinct}, for every $n \geq 0$ and
every pair of distinct vertices $v, w \in V_n$,
$\|x_w^{(n)} - x_v^{(n)}\| > 0$.
\end{corollary}

\begin{proof}
Immediate from Hypothesis~\ref{fact:distinct} and
positive-definiteness of the Euclidean norm on $\R^4$.
\end{proof}

\begin{notation}[$W(D_4)$ action]
\label{not:wd4}
Let $W(D_4) \leq O(4)$ denote the Coxeter group acting
orthogonally on $\R^4$ as the Weyl group of the root system
$D_4$ \cite[Ch.~1]{Humphreys}, realized concretely on the $D_4$
coordinate $4$-plane of \cite[\S 3]{AlgebraicSubstrate}.
\cite[Prop.~2.6]{RefinementSpaces} states that $W(D_4)$ permutes
$V_n$ for every $n$ via the Schl\"{a}fli-refinement-preserving
extension of its base action on $V_0 = V_{24}$. Coordinate
equivariance of this vertex action is proved below in
Lemma~\ref{lem:coord-equivariance}.
\end{notation}

\begin{lemma}[Coordinate equivariance on refined vertices]
\label{lem:coord-equivariance}
For every $n \geq 0$, every $g \in W(D_4)$, and every $v \in V_n$,
\[
  x_{g \cdot v}^{(n)} \;=\; g\,x_v^{(n)} \quad \text{in } \R^4.
\]
\end{lemma}

\begin{proof}
Induction on $n$.

\emph{Base case $n = 0$.} $V_0 = V_{24}$ is the $24$-cell vertex
set in the coordinate realization of
\cite[Def.~5.3]{AlgebraicSubstrate}, on which the Coxeter group
$W(D_4) \leq O(4)$ acts by its orthogonal action on $\R^4$
restricted to $V_{24}$ \cite[\S 3, Cor.~5.5]{AlgebraicSubstrate}.
With the coordinate identification $x_v^{(0)} = v \in \R^4$ for
$v \in V_{24}$, the claim $x_{g \cdot v}^{(0)} = g\,v = g\,x_v^{(0)}$
is the definition of the $W(D_4)$-action on $V_0$.

\emph{Inductive step.} Assume $x_{g \cdot u}^{(n)} = g\,x_u^{(n)}$
for all $g \in W(D_4)$ and $u \in V_n$. Level-$(n+1)$ vertices
are built from level-$n$ data as:
\begin{enumerate}
\item[\textup{(a)}] inclusion of old vertices
  $V_n \subset V_{n+1}$, with $x_u^{(n+1)} = x_u^{(n)}$
  (Notation~\ref{not:vertices-coords});
\item[\textup{(b)}] edge-midpoint vertices: for each unoriented
  level-$n$ edge $e = \{u_1, u_2\} \in E(G_n^{D_4})$, a new vertex
  $m_e$ with $x_{m_e}^{(n+1)} = \tfrac{1}{2}(x_{u_1}^{(n)} +
  x_{u_2}^{(n)})$, per \cite[Def.~2.2]{RefinementSpaces};
\item[\textup{(c)}] face-centroid vertices: for each level-$n$
  $2$-face $F \subset G_n^{D_4}$ with vertex set
  $\{u_1, \ldots, u_{k_F}\}$, a new vertex $c_F$ with
  $x_{c_F}^{(n+1)} = k_F^{-1}\sum_{j=1}^{k_F} x_{u_j}^{(n)}$,
  per \cite[Def.~2.2]{RefinementSpaces}.
\end{enumerate}
By \cite[Prop.~2.6]{RefinementSpaces}, the $W(D_4)$-action on
$V_{n+1}$ is the Schl\"{a}fli-refinement-preserving extension of
its action on $V_n$, and in particular:
$g \cdot u = g \cdot u$ for $u \in V_n$ (case (a), same as
level-$n$ action); $g \cdot m_{\{u_1, u_2\}} = m_{\{g u_1,
g u_2\}}$ (case (b), $W(D_4)$ permutes edges); and
$g \cdot c_F = c_{g \cdot F}$ (case (c), $W(D_4)$ permutes faces).

For case (a), $x_{g \cdot u}^{(n+1)} = x_{g \cdot u}^{(n)} =
g\,x_u^{(n)} = g\,x_u^{(n+1)}$ by inductive hypothesis.

For case (b), using the midpoint formula,
\[
  x_{g \cdot m_{\{u_1, u_2\}}}^{(n+1)}
    \;=\; x_{m_{\{g u_1, g u_2\}}}^{(n+1)}
    \;=\; \tfrac{1}{2}\bigl(x_{g u_1}^{(n)} + x_{g u_2}^{(n)}\bigr)
    \;=\; \tfrac{1}{2}\bigl(g\,x_{u_1}^{(n)} + g\,x_{u_2}^{(n)}\bigr)
    \;=\; g\,x_{m_{\{u_1, u_2\}}}^{(n+1)},
\]
using the inductive hypothesis in the third equality and
linearity of $g \in O(4)$ in the fourth.

For case (c), using the centroid formula,
\[
  x_{g \cdot c_F}^{(n+1)}
    \;=\; x_{c_{g \cdot F}}^{(n+1)}
    \;=\; k_F^{-1}\sum_{j=1}^{k_F} x_{g u_j}^{(n)}
    \;=\; k_F^{-1}\sum_{j=1}^{k_F} g\,x_{u_j}^{(n)}
    \;=\; g\,x_{c_F}^{(n+1)},
\]
using the inductive hypothesis and linearity of $g$. Induction
completes the proof.
\end{proof}

\begin{fact}[Vertex sets are finite]
\label{fact:finite}
$|V_n| < \infty$ for every $n \geq 0$.
\end{fact}

\begin{proof}
By induction from $|V_0| = 24$ and finiteness of the level-$n$
edge and $2$-face sets of the Schl\"{a}fli refinement
\cite[Prop.~2.6]{RefinementSpaces}.
\end{proof}

\begin{notation}[Coordinate diameter]
\label{not:diam}
Set
\[
  d_n \;:=\; \min_{v \neq w \in V_n}
    \|x_w^{(n)} - x_v^{(n)}\|, \qquad
  D_n \;:=\; \max_{v, w \in V_n}
    \|x_w^{(n)} - x_v^{(n)}\|.
\]
Under Hypothesis~\ref{fact:distinct} (and therefore
Corollary~\ref{cor:nonvanishing}) we have $0 < d_n \leq D_n <
\infty$, and $D_n$ is bounded by the $24$-cell diameter, hence by
a finite constant independent of $n$.
\end{notation}

%=====================================================================
\section{Source Data from the P4 State Space}
\label{sec:source}
%=====================================================================

\begin{definition}[Scalar source space]
\label{def:scalar-source}
The \emph{scalar source space} at level $n$ is
\[
  \mathcal S_n \;:=\; \ell^2(V_n),
\]
the finite-dimensional Hilbert space of $\R$-valued functions on
$V_n$ with the counting-measure inner product
$\langle s, s' \rangle := \sum_{v \in V_n} s(v)\,s'(v)$.
\end{definition}

\begin{definition}[P4 state space on the $D_4$ branch]
\label{def:p4-state}
Let $X_n := X_n^{D_4} = X_n^{0, D_4} \oplus X_n^{1, D_4}$ denote
the finite-level P4 closure-dynamics state space
\cite[Def.~2.1]{ClosureDynamics}, with the direct-sum inner
product.
\end{definition}

\begin{definition}[Admissible fibre functional]
\label{def:fibre-functional}
A \emph{$W(D_4)$-invariant fibre functional} is a bounded
$\R$-linear map $\lambda \colon F^0 \to \R$ satisfying
\[
  \lambda \circ \rho_{F^0}^{D_4}(g) \;=\; \lambda
  \qquad \text{for every } g \in W(D_4),
\]
where $F^0$ is the $D_4$-branch vertex fibre and
$\rho_{F^0}^{D_4} \colon W(D_4) \to O(F^0)$ is the branchwise
Coxeter action of \cite[Def.~8.1]{RefinementSpaces}. The
operator norm $\|\lambda\|_\op$ is the norm of $\lambda$ as an
element of the finite-dimensional dual $(F^0)^\ast$.
\end{definition}

\begin{proposition}[Existence of nonzero invariant fibre functionals]
\label{prop:lambda-exists}
There exists at least one nonzero $W(D_4)$-invariant fibre
functional $\lambda \in (F^0)^\ast$.
\end{proposition}

\begin{proof}
By \cite[Def.~8.1]{RefinementSpaces} (branchwise Coxeter action),
on the $D_4$ branch the fibre representation
$\rho_{F^0}^{D_4} \colon W(D_4) \to O(F^0)$ acts by the standard
$D_4$-reflection representation on the $V_{D_4}$ summand of
$F^0$ and \emph{trivially} on the $V_{H_4}, V_{16}, V_8$
summands. Therefore the invariant subspace satisfies
\[
  (F^0)^{W(D_4)} \;\supseteq\; V_{H_4} \oplus V_{16} \oplus V_8,
\]
and in particular $(F^0)^{W(D_4)}$ is nonzero. Choose any nonzero
vector $w_0 \in V_{16}$ (say) and any norm-one dual basis element
$\varphi_0 \in V_{16}^\ast$ with $\varphi_0(w_0) \neq 0$; extend
$\varphi_0$ to a linear functional $\lambda \in (F^0)^\ast$ by
setting $\lambda(v) := \varphi_0(\mathrm{pr}_{V_{16}} v)$, where
$\mathrm{pr}_{V_{16}} \colon F^0 \to V_{16}$ is the direct-sum
projection. Then $\lambda$ vanishes on $V_{H_4} \oplus V_{D_4}
\oplus V_8$ and equals $\varphi_0$ on $V_{16}$. For every
$g \in W(D_4)$, the action $\rho_{F^0}^{D_4}(g)$ preserves the
$V_{16}$ summand and acts trivially on it
\cite[Def.~8.1]{RefinementSpaces}, so
$\lambda(\rho_{F^0}^{D_4}(g) v) = \varphi_0(\mathrm{pr}_{V_{16}}
(\rho_{F^0}^{D_4}(g) v)) = \varphi_0(\mathrm{pr}_{V_{16}} v) =
\lambda(v)$. Hence $\lambda \circ \rho_{F^0}^{D_4}(g) = \lambda$
for every $g$, and $\lambda(w_0) = \varphi_0(w_0) \neq 0$ shows
$\lambda$ is nonzero. Its operator norm is
$\|\lambda\|_\op = \|\varphi_0\|_\op$, finite by construction.
\end{proof}

\begin{remark}[Averaging construction of further invariant functionals]
\label{rmk:lambda-averaging}
An alternative construction that works for any finite group
representation is the group average. For any
$\varphi \in (F^0)^\ast$, set
\[
  \lambda_\varphi \;:=\; \frac{1}{|W(D_4)|}
    \sum_{g \in W(D_4)}
    \varphi \circ \rho_{F^0}^{D_4}(g).
\]
$\lambda_\varphi$ is $W(D_4)$-invariant by the standard
substitution-of-summation argument; it equals the orthogonal
projection of $\varphi$ onto the $W(D_4)$-invariant summand of
$(F^0)^\ast$. The norm bound $\|\lambda_\varphi\|_\op \leq
\|\varphi\|_\op$ follows from the triangle inequality applied to
the averaged sum and the orthogonality of each
$\rho_{F^0}^{D_4}(g)$ on $F^0$
(which implies $\|\varphi \circ \rho_{F^0}^{D_4}(g)\|_\op =
\|\varphi\|_\op$; this orthogonality is printed in
\cite[Def.~8.1]{RefinementSpaces} as part of the action's
definition). Proposition~\ref{prop:lambda-exists} gives a
concrete nonzero instance; the averaging remark records the
general construction.
\end{remark}

\begin{definition}[Source extraction map $S_n^\lambda$]
\label{def:source-extraction}
Fix a $W(D_4)$-invariant fibre functional $\lambda \in (F^0)^\ast$
(Definition~\ref{def:fibre-functional}), chosen once and for all.
The \emph{source extraction operator}
$S_n^\lambda \colon X_n \to \mathcal S_n$ is the composition of:
\begin{itemize}
\item the direct-sum vertex-sector projection
  $\pi^0 \colon X_n^{0, D_4} \oplus X_n^{1, D_4} \to X_n^{0, D_4}$,
  and
\item the $\lambda$-pullback $\tau_\lambda \colon
  X_n^{0, D_4} \to \mathcal S_n$ defined at each vertex
  $v \in V_n$ by
  \[
    \tau_\lambda(f)(v) \;:=\; \lambda(f(v)),
  \]
  where $f \in X_n^{0, D_4}$ is the standard fibrewise
  representation $f = (f(v))_{v \in V_n}$ with
  $f(v) \in F^0$ (the P3 sector fibre
  \cite[Notation~3.3]{RefinementSpaces}).
\end{itemize}
Thus $S_n^\lambda := \tau_\lambda \circ \pi^0$.
\end{definition}

\begin{proposition}[Source extraction is bounded]
\label{prop:source-bounded}
For every $W(D_4)$-invariant fibre functional $\lambda$, the
operator $S_n^\lambda \colon X_n \to \mathcal S_n$ is linear and
bounded with
$\|S_n^\lambda\|_\op \leq \|\lambda\|_\op$.
\end{proposition}

\begin{proof}
$S_n^\lambda$ is linear by linearity of $\lambda$ and of $\pi^0$.
For each $v \in V_n$, $|\lambda(f(v))| \leq \|\lambda\|_\op
\cdot \|f(v)\|_{F^0}$ by the definition of the dual norm. Squaring
and summing over $v \in V_n$,
$\|\tau_\lambda(f)\|_{\mathcal S_n}^2 = \sum_v |\lambda(f(v))|^2
\leq \|\lambda\|_\op^2 \sum_v \|f(v)\|_{F^0}^2 =
\|\lambda\|_\op^2 \|f\|_{X_n^0}^2$,
so $\|\tau_\lambda\|_\op \leq \|\lambda\|_\op$. Since $\pi^0$ is
the orthogonal direct-sum projection with $\|\pi^0\|_\op \leq 1$,
composition gives
$\|S_n^\lambda\|_\op \leq \|\tau_\lambda\|_\op \cdot \|\pi^0\|_\op
\leq \|\lambda\|_\op$.
\end{proof}

\begin{remark}[Choice of $\lambda$]
\label{rmk:source-choice}
Every theorem below is stated for a fixed but arbitrary
$W(D_4)$-invariant fibre functional $\lambda$. The metric
projection $\Pimet^{(n)} = P_n^\PF \circ M_n \circ S_n^\lambda$
inherits $\lambda$ as a parameter, and every operator-norm bound
involving $\Pimet^{(n)}$ carries $\|\lambda\|_\op$. This choice
keeps P5 independent of any unprinted fibre-level operator
structure beyond what P3 already prints (the finite-dimensional
real vector space $F^0$ and the unitary Coxeter action
$\rho_{F^0}^{D_4}$). In downstream work, any specific
computation of $\Pimet^{(n)}$ requires first fixing a specific
$\lambda$.
\end{remark}

%=====================================================================
\section{The Discrete Moment Tensor}
\label{sec:moment}
%=====================================================================

\begin{definition}[Normalized rank-one displacement kernel]
\label{def:kernel}
For every $n \geq 0$ and every pair of distinct vertices
$v, w \in V_n$, set
\[
  K_n(v, w) \;:=\;
    \frac{(x_w^{(n)} - x_v^{(n)}) \otimes
      (x_w^{(n)} - x_v^{(n)})}{\|x_w^{(n)} - x_v^{(n)}\|^2}
  \;\in\; \Sym^2(\R^4).
\]
Here $(a \otimes a)$ is the rank-one symmetric matrix with
entries $a_i a_j$. By Corollary~\ref{cor:nonvanishing} the
denominator is nonzero, so $K_n(v, w)$ is well-defined.
\end{definition}

\begin{lemma}[Properties of $K_n$]
\label{lem:kernel-properties}
For all distinct $v, w \in V_n$:
\begin{enumerate}
\item[\textup{(i)}] $K_n(v, w) \in \Sym^2(\R^4)$, i.e.\
  $K_n(v, w)^\top = K_n(v, w)$.
\item[\textup{(ii)}] $K_n(v, w)$ is a rank-one orthogonal
  projector: $K_n(v, w)^2 = K_n(v, w)$ and
  $\tr K_n(v, w) = 1$.
\item[\textup{(iii)}] $K_n(v, w) = K_n(w, v)$.
\item[\textup{(iv)}] $\|K_n(v, w)\|_F = 1$ in the Frobenius
  norm, equivalently $\|K_n(v, w)\|_\op = 1$ as an operator on
  $\R^4$.
\end{enumerate}
\end{lemma}

\begin{proof}
Write $r := x_w^{(n)} - x_v^{(n)}$ and $\hat r := r / \|r\|$.
Then $K_n(v, w) = \hat r \otimes \hat r = \hat r\,\hat r^\top$.
(i): $(a \otimes a)^\top = a \otimes a$. (ii): for any unit
vector $\hat r$, $(\hat r\,\hat r^\top)^2 = \hat r\,(\hat r^\top
\hat r)\,\hat r^\top = \hat r\,\hat r^\top$, and $\tr(\hat r\,
\hat r^\top) = \hat r^\top \hat r = 1$. (iii):
$(x_v - x_w) = -(x_w - x_v)$ and $(-a) \otimes (-a) = a \otimes
a$. (iv): the Frobenius norm of a rank-one projector onto a unit
direction is $\sqrt{\tr(K^2)} = \sqrt{\tr K} = 1$; the operator
norm equals the largest eigenvalue, which is $1$.
\end{proof}

\begin{definition}[Raw symmetric-tensor space]
\label{def:tn}
Let
\[
  \mathcal T_n \;:=\; \ell^2(V_n; \Sym^2(\R^4))
\]
be the finite-dimensional Hilbert space of $\Sym^2(\R^4)$-valued
functions on $V_n$ with the counting-measure Frobenius inner
product
\[
  \langle T, T' \rangle_{\mathcal T_n}
    \;:=\; \sum_{v \in V_n} \tr\bigl(T(v)^\top T'(v)\bigr)
    \;=\; \sum_{v \in V_n} \langle T(v), T'(v)
    \rangle_F.
\]
\end{definition}

\begin{definition}[Discrete moment tensor operator]
\label{def:mn}
The \emph{discrete moment tensor operator} at level $n$ is the
linear map $M_n \colon \mathcal S_n \to \mathcal T_n$ defined
pointwise by
\[
  (M_n s)(v) \;:=\; \sum_{w \in V_n,\ w \neq v} s(w)\,K_n(v, w),
  \qquad v \in V_n.
\]
The sum has finitely many terms by Fact~\ref{fact:finite}.
\end{definition}

\begin{theorem}[$M_n$ is well-defined, bounded, and symmetric]
\label{thm:mn-welldef}
For every $n \geq 0$:
\begin{enumerate}
\item[\textup{(i)}] The operator $M_n \colon \mathcal S_n \to
  \mathcal T_n$ is a well-defined linear map, and for every
  $s \in \mathcal S_n$ and every $v \in V_n$, $(M_n s)(v) \in
  \Sym^2(\R^4)$.
\item[\textup{(ii)}] $M_n$ is bounded, with operator-norm bound
  $\|M_n\|_\op \leq |V_n|$.
\item[\textup{(iii)}] For pointwise nonnegative $s \geq 0$ and
  every $v$, $(M_n s)(v)$ is a positive-semidefinite symmetric
  matrix.
\end{enumerate}
\end{theorem}

\begin{proof}
(i): By Corollary~\ref{cor:nonvanishing} and
Fact~\ref{fact:finite}, the defining sum in
Definition~\ref{def:mn} is a finite sum of well-defined symmetric
rank-one matrices (Lemma~\ref{lem:kernel-properties}(i)). The sum
of symmetric matrices is symmetric, so $(M_n s)(v) \in
\Sym^2(\R^4)$. Linearity in $s$ is immediate from the defining
formula.

(ii): For each $v \in V_n$,
\[
  \|(M_n s)(v)\|_F
    \;=\; \biggl\|\sum_{w \neq v} s(w)\,K_n(v, w)\biggr\|_F
    \;\leq\; \sum_{w \neq v} |s(w)|\,\|K_n(v, w)\|_F
    \;=\; \sum_{w \neq v} |s(w)|
\]
by Lemma~\ref{lem:kernel-properties}(iv). By the Cauchy--Schwarz
inequality on $V_n \setminus \{v\}$,
\[
  \sum_{w \neq v} |s(w)|
    \;\leq\; \sqrt{|V_n| - 1} \cdot \|s\|_{\mathcal S_n}
    \;\leq\; \sqrt{|V_n|} \cdot \|s\|_{\mathcal S_n}.
\]
Squaring and summing over $v \in V_n$,
\[
  \|M_n s\|_{\mathcal T_n}^2
    \;=\; \sum_v \|(M_n s)(v)\|_F^2
    \;\leq\; \sum_v |V_n|\cdot \|s\|_{\mathcal S_n}^2
    \;=\; |V_n|^2\,\|s\|_{\mathcal S_n}^2,
\]
whence $\|M_n s\|_{\mathcal T_n} \leq |V_n|\cdot
\|s\|_{\mathcal S_n}$, giving $\|M_n\|_\op \leq |V_n|$.

(iii): For $s \geq 0$, $(M_n s)(v) = \sum_{w \neq v}
s(w)\,K_n(v, w)$ is a nonnegative combination of rank-one
projectors $K_n(v, w)$. Each $K_n(v, w)$ is positive
semidefinite by Lemma~\ref{lem:kernel-properties}(ii), and the
cone of positive semidefinite symmetric matrices is closed under
nonnegative linear combinations. Hence $(M_n s)(v) \succeq 0$.
\end{proof}

\begin{remark}[Level-dependent bound]
\label{rmk:level-bound}
The bound $\|M_n\|_\op \leq |V_n|$ is level-dependent and grows
polynomially with $n$ via $|V_n|$. No finer asymptotic is
claimed here. Convergence under refinement of the operator
family $\{M_n\}_{n \geq 0}$ is the target of P6 and is out of
scope for this paper.
\end{remark}

%=====================================================================
\section{The Piecewise-Flat Target and the Metric Projection}
\label{sec:pf-target}
%=====================================================================

\subsection{Cell-labelled symmetric tensor fields}

The piecewise-flat target in this paper is the space of
\emph{cell-labelled} symmetric tensor fields at level $n$: each
vertex $v \in V_n$ plays the role of an abstract cell label, and
each cell carries a constant symmetric tensor value in
$\Sym^2(\R^4)$. No cross-cell (``face-compatibility'') constraint
between distinct cells is imposed in this model, and no global
continuous field on a geometric realization of the refinement
complex is constructed. The name ``piecewise-flat'' indicates
only that the tensor is constant on each cell; it does not imply
a continuous or geometrically-realized field. A future paper
that wishes to impose shared-face compatibility conditions for a
genuinely piecewise-flat global field may do so by passing to a
named subspace of $\PF_n$ with a printed compatibility rule; no
such additional structure is used here.

\begin{definition}[Piecewise-flat target $\PF_n$]
\label{def:pfn}
The \emph{piecewise-flat target} at level $n$ is the
finite-dimensional real Hilbert space
\[
  \PF_n \;:=\; \ell^2(V_n;\Sym^2(\R^4)),
\]
the space of functions $T \colon V_n \to \Sym^2(\R^4)$, with the
\emph{cellwise Frobenius inner product}
\[
  \langle T, T' \rangle_{\PF_n}
    \;:=\; \sum_{v \in V_n} \langle T(v), T'(v) \rangle_F,
  \qquad \langle A, B \rangle_F := \tr(A^\top B),
\]
and the induced norm
$\|T\|_{\PF_n}^2 = \sum_v \|T(v)\|_F^2$. The admissibility
condition on a cell-labelled tensor field $T \in \PF_n$ is: for
every $v \in V_n$, $T(v) \in \Sym^2(\R^4)$ (symmetry on each
cell). No cross-cell compatibility condition is imposed.
\end{definition}

\begin{remark}[$\PF_n$ vs.\ $\mathcal T_n$]
\label{rmk:pfn-vs-tn}
As finite-dimensional real Hilbert spaces, $\PF_n$ and
$\mathcal T_n$ are literally identical: both are
$\ell^2(V_n;\Sym^2(\R^4))$ with the cellwise Frobenius inner
product. The notational distinction is purely interpretive:
$\mathcal T_n$ is the raw codomain of $M_n$
(Definition~\ref{def:tn}); $\PF_n$ is the same space viewed as
the cell-labelled target for $\Pimet^{(n)}$. The admissibility
condition of Definition~\ref{def:pfn} (symmetry of each $T(v)$)
is already built into the codomain $\Sym^2(\R^4)$ of $M_n$, so
no nontrivial projection is required to pass from $\mathcal T_n$
to $\PF_n$. This is why the piecewise-flat projection
(Definition~\ref{def:pn-pf}) is the identity in this model.
\end{remark}

\begin{definition}[Piecewise-flat projection $P_n^\PF$]
\label{def:pn-pf}
The \emph{piecewise-flat projection} is the canonical identity
map
\[
  P_n^\PF \;:=\; \id_{\mathcal T_n}
    \;\colon\; \mathcal T_n \to \PF_n.
\]
It is an isometric isomorphism of Hilbert spaces, in particular
$\|P_n^\PF\|_\op = 1$.
\end{definition}

\begin{remark}[Why $P_n^\PF = \id$ here]
\label{rmk:pnpf-identity}
The only admissibility condition imposed by
Definition~\ref{def:pfn} is symmetry of each cell tensor
$T(v) \in \Sym^2(\R^4)$, which is automatic because the codomain
of $M_n$ is already $\Sym^2(\R^4)$. No cross-cell compatibility
condition is imposed (Remark~\ref{rmk:pfn-vs-tn}). No further
projection is required, so $P_n^\PF$ is the identity. We record it as a named bounded
linear operator so that every theorem involving $\Pimet^{(n)}$
can be stated as a composition of three named bounded linear
maps, in the standard form of the work order. Any later paper
that imposes additional admissibility constraints (e.g.\ a
codimension-$1$ tangential-space constraint, or a trace
normalization) would redefine $P_n^\PF$ as the corresponding
orthogonal projection; all theorems below would then carry the
reduced bound $\|P_n^\PF\|_\op \leq 1$ instead of equality.
\end{remark}

\begin{definition}[Metric projection $\Pimet^{(n)}$]
\label{def:pimet}
The \emph{level-$n$ metric projection} is the composition
\[
  \Pimet^{(n)} \;:=\; P_n^\PF \circ M_n \circ S_n^\lambda
    \;\colon\; X_n \to \PF_n.
\]
When we wish to consider the source-level variant we write
$\widetilde{\Pimet}^{(n)} := P_n^\PF \circ M_n \colon
\mathcal S_n \to \PF_n$, so that $\Pimet^{(n)} =
\widetilde{\Pimet}^{(n)} \circ S_n^\lambda$.
\end{definition}

%=====================================================================
\section{Refinement Stability}
\label{sec:refinement}
%=====================================================================

\begin{notation}[Refinement bonding maps]
\label{not:bonding}
Let $p_{n+1, n}^0 \colon X_{n+1}^{0, D_4} \to X_n^{0, D_4}$ be
the P3 vertex bonding by restriction
\cite[Def.~5.1]{RefinementSpaces}; extend componentwise to
$p_{n+1, n}^0 \colon \mathcal S_{n+1} \to \mathcal S_n$ via
\[
  (p_{n+1, n}^0 s)(v) \;:=\; s(v), \qquad v \in V_n \subset V_{n+1},
\]
using the vertex inclusion of \cite[Rmk.~2.3]{RefinementSpaces}.
Let $p_{n+1, n}^X \colon X_{n+1} \to X_n$ be the direct-sum
bonding $p_{n+1, n}^X := p_{n+1, n}^0 \oplus p_{n+1, n}^1$ of
\cite[Thm.~5.3]{RefinementSpaces}. Both are contractions.
\end{notation}

\begin{definition}[Tensor coarse projection]
\label{def:qn}
The \emph{tensor coarse projection} $q_{n+1, n} \colon
\mathcal T_{n+1} \to \mathcal T_n$ is defined by vertex
restriction:
\[
  (q_{n+1, n} T)(v) \;:=\; T(v), \qquad v \in V_n \subset V_{n+1}.
\]
The piecewise-flat analogue $q_{n+1, n}^\PF \colon \PF_{n+1} \to
\PF_n$ is the same formula under the identification $\PF_n =
\mathcal T_n$.
\end{definition}

\begin{lemma}[Source--state bonding compatibility]
\label{lem:source-state-bonding}
For every $W(D_4)$-invariant fibre functional $\lambda$, the
source extraction satisfies
\[
  S_n^\lambda \circ p_{n+1, n}^X
    \;=\; p_{n+1, n}^0 \circ S_{n+1}^\lambda
  \qquad \text{on } X_{n+1} \to \mathcal S_n.
\]
\end{lemma}

\begin{proof}
By Definition~\ref{def:source-extraction}, $S_n^\lambda =
\tau_\lambda \circ \pi^0$ with $\pi^0$ the direct-sum
vertex-sector projection. The direct-sum bonding $p_{n+1, n}^X =
p_{n+1, n}^0 \oplus p_{n+1, n}^1$ commutes with $\pi^0$:
$\pi^0 \circ p_{n+1, n}^X = p_{n+1, n}^0 \circ \pi^0$, because
$\pi^0$ is the projection onto the first summand and
$p_{n+1, n}^X$ preserves summands. The $\lambda$-pullback
$\tau_\lambda$ commutes with vertex restriction of
fibre-valued functions: for $v \in V_n \subset V_{n+1}$ and
$f \in X_{n+1}^{0, D_4}$,
$(\tau_\lambda \circ p_{n+1, n}^0)(f)(v) =
\lambda((p_{n+1, n}^0 f)(v)) = \lambda(f(v)) =
\tau_\lambda(f)(v) = (p_{n+1, n}^0 \circ \tau_\lambda)(f)(v)$,
using the vertex inclusion $V_n \subset V_{n+1}$ and the formula
$(p_{n+1, n}^0 f)(v) = f(v)$ from Notation~\ref{not:bonding}.
Composing:
$S_n^\lambda \circ p_{n+1, n}^X = \tau_\lambda \circ \pi^0
\circ p_{n+1, n}^X = \tau_\lambda \circ p_{n+1, n}^0 \circ \pi^0 =
p_{n+1, n}^0 \circ \tau_\lambda \circ \pi^0 = p_{n+1, n}^0 \circ
S_{n+1}^\lambda$.
\end{proof}

\begin{definition}[Old-vertex subspace]
\label{def:old-vertex-subspace}
The \emph{old-vertex subspace} at level $n+1$ is
\[
  \mathcal S_{n+1}^{\mathrm{old}}
    \;:=\; \{\,s \in \mathcal S_{n+1} :
      s(w) = 0 \text{ for every } w \in V_{n+1}\setminus V_n\,\},
\]
the finite-dimensional subspace of level-$(n+1)$ source functions
supported on the old-vertex set $V_n \subset V_{n+1}$. Under the
inclusion-induced zero-extension map
$\iota_{n, n+1} \colon \mathcal S_n \to \mathcal S_{n+1}^{\mathrm{old}}$,
$(\iota_{n, n+1} s)(w) := s(w)$ if $w \in V_n$ and $0$ otherwise,
$\iota_{n, n+1}$ is a linear isometric bijection and
$p_{n+1, n}^0 \circ \iota_{n, n+1} = \id_{\mathcal S_n}$.

Analogously on the state side,
\[
  X_{n+1}^{\mathrm{old}} \;:=\;
    (S_{n+1}^\lambda)^{-1}\bigl(\mathcal S_{n+1}^{\mathrm{old}}\bigr)
      \;\subset\; X_{n+1}
\]
is the preimage of $\mathcal S_{n+1}^{\mathrm{old}}$ under the
level-$(n+1)$ source extraction.
\end{definition}

\begin{theorem}[Refinement stability of $M_n$ on old-vertex sources]
\label{thm:mn-refinement}
For every $n \geq 0$, the operator identity
\[
  q_{n+1, n} \circ M_{n+1}|_{\mathcal S_{n+1}^{\mathrm{old}}}
    \;=\; M_n \circ p_{n+1, n}^0|_{\mathcal S_{n+1}^{\mathrm{old}}}
  \qquad \text{as maps from }
    \mathcal S_{n+1}^{\mathrm{old}} \text{ to } \mathcal T_n
\]
holds. Consequently, on the state side,
\[
  q_{n+1, n}^\PF \circ
    \Pimet^{(n+1)}|_{X_{n+1}^{\mathrm{old}}}
    \;=\; \Pimet^{(n)} \circ
       p_{n+1, n}^X|_{X_{n+1}^{\mathrm{old}}}
  \qquad \text{as maps from }
    X_{n+1}^{\mathrm{old}} \text{ to } \PF_n.
\]
The identity is stated only on the old-vertex subspaces; it does
not extend to all of $\mathcal S_{n+1}$ or $X_{n+1}$ because the
level-$(n+1)$ source at new vertices $V_{n+1}\setminus V_n$
produces additional nonzero kernel terms in $M_{n+1}s(v)$ that are
not captured by $M_n p_{n+1, n}^0 s$ (see
Remark~\ref{rmk:refinement-scope}).
\end{theorem}

\begin{proof}
Fix $s \in \mathcal S_{n+1}^{\mathrm{old}}$ and $v \in V_n
\subset V_{n+1}$. By the defining condition, $s(w) = 0$ for every
$w \in V_{n+1}\setminus V_n$. On the one hand, by Definitions
\ref{def:mn} and \ref{def:qn},
\[
  (q_{n+1, n} M_{n+1} s)(v)
    \;=\; (M_{n+1} s)(v)
    \;=\; \sum_{w \in V_{n+1},\ w \neq v} s(w)\,K_{n+1}(v, w)
    \;=\; \sum_{w \in V_n\setminus\{v\}} s(w)\,K_{n+1}(v, w),
\]
using $s(w) = 0$ for $w \in V_{n+1}\setminus V_n$ to drop the
new-vertex summands. By Notation~\ref{not:vertices-coords},
$x_u^{(n+1)} = x_u^{(n)}$ for every $u \in V_n$, so
$K_{n+1}(v, w) = K_n(v, w)$ whenever $v, w \in V_n$. Hence
\[
  (q_{n+1, n} M_{n+1} s)(v)
    \;=\; \sum_{w \in V_n\setminus\{v\}} s(w)\,K_n(v, w).
\]

On the other hand, by Definition~\ref{def:mn} and
Notation~\ref{not:bonding},
\[
  (M_n p_{n+1, n}^0 s)(v)
    \;=\; \sum_{w \in V_n\setminus\{v\}}
       (p_{n+1, n}^0 s)(w)\,K_n(v, w)
    \;=\; \sum_{w \in V_n\setminus\{v\}} s(w)\,K_n(v, w),
\]
using $(p_{n+1, n}^0 s)(w) = s(w)$ for $w \in V_n$. The two sides
agree, proving the source-level identity on
$\mathcal S_{n+1}^{\mathrm{old}}$.

The state-level consequence follows by composing with
$P_n^\PF = \id$ and
Lemma~\ref{lem:source-state-bonding}: for every
$\Phi \in X_{n+1}^{\mathrm{old}}$,
$S_{n+1}^\lambda \Phi \in \mathcal S_{n+1}^{\mathrm{old}}$ by
definition of $X_{n+1}^{\mathrm{old}}$, and
\begin{align*}
  (q_{n+1, n}^\PF \Pimet^{(n+1)} \Phi)
    &\;=\; q_{n+1, n}^\PF\,P_{n+1}^\PF\,M_{n+1}\,S_{n+1}^\lambda
       \,\Phi
     \;=\; q_{n+1, n}\,M_{n+1}\,S_{n+1}^\lambda\,\Phi \\
    &\;=\; M_n\,p_{n+1, n}^0\,S_{n+1}^\lambda\,\Phi
     \;=\; M_n\,S_n^\lambda\,p_{n+1, n}^X\,\Phi
     \;=\; \Pimet^{(n)}\,p_{n+1, n}^X\,\Phi,
\end{align*}
using $P_{n+1}^\PF = \id$, $P_n^\PF = \id$, the source-level
identity just proved applied to
$S_{n+1}^\lambda \Phi \in \mathcal S_{n+1}^{\mathrm{old}}$, and
Lemma~\ref{lem:source-state-bonding}.
\end{proof}

\begin{remark}[Honest scope of refinement stability]
\label{rmk:refinement-scope}
Theorem~\ref{thm:mn-refinement}, as proved, holds as a literal
operator identity only on the subspace of level-$(n+1)$ source
functions vanishing on newly added vertices. This is the honest
finite-level statement: the moment tensor is determined by the
source values at all visible vertices, and refinement adds new
vertices that carry their own source data. The identity cannot be
extended to all of $\mathcal S_{n+1}$ without introducing a
projection or discarding the new-vertex contributions, and we
decline to do either. Convergence of the full $M_{n+1}$ to some
limit of $M_n$ as $n \to \infty$ is outside the scope of this
paper; that question is the target of P6.
\end{remark}

%=====================================================================
\section{$W(D_4)$-Equivariance}
\label{sec:equivariance}
%=====================================================================

\begin{notation}[Induced $W(D_4)$ actions]
\label{not:induced-actions}
The orthogonal $W(D_4)$ action on $\R^4$ of
Notation~\ref{not:wd4} induces:
\begin{itemize}
\item on $\mathcal S_n = \ell^2(V_n)$, the permutation action
  $(g \cdot s)(v) := s(g^{-1} v)$;
\item on $\mathcal T_n = \ell^2(V_n; \Sym^2(\R^4))$, the twisted
  permutation action
  $(g \cdot T)(v) := g\,T(g^{-1} v)\,g^\top$, where the
  right-hand side is the conjugation action of $g \in O(4)$ on
  $\Sym^2(\R^4)$;
\item on $\PF_n = \mathcal T_n$, the same formula;
\item on $X_n$, the P3 branchwise Coxeter action
  \cite[Def.~8.1]{RefinementSpaces}.
\end{itemize}
\end{notation}

\begin{fact}[All induced actions are orthogonal]
\label{fact:actions-orthogonal}
Each of the four induced actions in
Notation~\ref{not:induced-actions} is a unitary representation
of $W(D_4)$ on the corresponding finite-dimensional Hilbert
space.
\end{fact}

\begin{proof}
For the $\mathcal S_n$ action this is the standard statement that
permutation actions preserve the counting-measure $\ell^2$ norm.
For the $\mathcal T_n$ and $\PF_n$ actions, $(g \cdot T)(v) =
g\,T(g^{-1}v)\,g^\top$ has Frobenius norm $\|g\,T(g^{-1}v)\,
g^\top\|_F = \|T(g^{-1}v)\|_F$, using that $g \in O(4)$ preserves
the Frobenius norm under conjugation (because conjugation by an
orthogonal matrix is an isometry on symmetric matrices in the
Frobenius inner product). Summing squared fibre norms over
$v \in V_n$ and using that $g$ permutes $V_n$ gives unitarity.
The $X_n$ case is \cite[Def.~8.1]{RefinementSpaces}, where the
unitarity of $\rho_{F^0}^{D_4}(g)$ on $X_n^{0, D_4}$ is stated as
part of the definition of the branchwise Coxeter action.
\end{proof}

\begin{lemma}[Kernel equivariance]
\label{lem:kernel-equiv}
For every $g \in W(D_4)$ and every distinct $v, w \in V_n$,
\[
  K_n(g v, g w) \;=\; g\,K_n(v, w)\,g^\top.
\]
\end{lemma}

\begin{proof}
By Lemma~\ref{lem:coord-equivariance}, $x_{g u}^{(n)} = g\,x_u^{(n)}$ for
every $u \in V_n$. Hence
$x_{g w}^{(n)} - x_{g v}^{(n)} = g(x_w^{(n)} - x_v^{(n)})$, and
$\|x_{g w}^{(n)} - x_{g v}^{(n)}\| = \|g(x_w^{(n)} - x_v^{(n)})\|
= \|x_w^{(n)} - x_v^{(n)}\|$ because $g \in O(4)$ preserves
Euclidean norm. Writing $r := x_w^{(n)} - x_v^{(n)}$ and
$\hat r := r/\|r\|$:
\[
  K_n(gv, gw)
    \;=\; \frac{(gr) \otimes (gr)}{\|gr\|^2}
    \;=\; \frac{(gr)(gr)^\top}{\|r\|^2}
    \;=\; g\,\frac{r\,r^\top}{\|r\|^2}\,g^\top
    \;=\; g\,K_n(v, w)\,g^\top.
\]
\end{proof}

\begin{theorem}[$W(D_4)$-equivariance of $M_n$ and $\Pimet^{(n)}$]
\label{thm:mn-equivariance}
For every $g \in W(D_4)$:
\begin{enumerate}
\item[\textup{(i)}] The moment operator intertwines the induced
  actions: $M_n(g \cdot s) = g \cdot M_n(s)$ for every $s \in
  \mathcal S_n$. Equivalently,
  $(M_n s)(gv) = g\,(M_n(g^{-1}\cdot s))(v)\,g^\top$ for every
  $v \in V_n$.
\item[\textup{(ii)}] The source extraction intertwines the
  actions: $S_n^\lambda(g \cdot \Phi) = g \cdot S_n^\lambda(\Phi)$ for every
  $\Phi \in X_n$.
\item[\textup{(iii)}] The piecewise-flat projection
  $P_n^\PF = \id$ intertwines the $\mathcal T_n$ and $\PF_n$
  actions (which coincide).
\item[\textup{(iv)}] Consequently, $\Pimet^{(n)}(g \cdot \Phi) =
  g \cdot \Pimet^{(n)}(\Phi)$ for every $\Phi \in X_n$.
\end{enumerate}
\end{theorem}

\begin{proof}
(i): Fix $g \in W(D_4)$, $s \in \mathcal S_n$, $v \in V_n$.
Substituting $(g \cdot s)(w) = s(g^{-1}w)$ and
$K_n(v, w) = g^{-1}\,K_n(gv, gw)\,(g^{-1})^\top$
(from Lemma~\ref{lem:kernel-equiv}, rewritten with $g^{-1}$), we
compute
\begin{align*}
  (M_n(g \cdot s))(v)
    &\;=\; \sum_{w \neq v} (g \cdot s)(w)\,K_n(v, w)
    \;=\; \sum_{w \neq v} s(g^{-1}w)\,K_n(v, w).
\end{align*}
Reparametrize $u := g^{-1}w$, so $w = gu$ and the sum over
$w \neq v$ becomes a sum over $u \neq g^{-1}v$:
\[
  (M_n(g \cdot s))(v)
    \;=\; \sum_{u \neq g^{-1}v} s(u)\,K_n(v, gu)
    \;=\; \sum_{u \neq g^{-1}v} s(u)\,g\,K_n(g^{-1}v, u)\,g^\top,
\]
using Lemma~\ref{lem:kernel-equiv} in the form $K_n(v, gu) =
g\,K_n(g^{-1}v, u)\,g^\top$. Factor $g$ and $g^\top$ outside the
sum:
\[
  (M_n(g \cdot s))(v)
    \;=\; g\,\biggl(\sum_{u \neq g^{-1}v}
       s(u)\,K_n(g^{-1}v, u)\biggr)\,g^\top
    \;=\; g\,(M_n s)(g^{-1}v)\,g^\top
    \;=\; (g \cdot M_n s)(v).
\]

(ii): For $\Phi = (f, B) \in X_n^{0, D_4} \oplus X_n^{1, D_4}$
and $g \in W(D_4)$, the P3 action
\cite[Def.~8.1]{RefinementSpaces} gives
$g \cdot (f, B) = (g \cdot f, g \cdot B)$ where
$(g \cdot f)(v) = \rho_{F^0}^{D_4}(g)\,f(g^{-1}v)$ in the P3
fibre $F^0$. The direct-sum projection $\pi^0$ commutes with the
direct-sum action because each summand is preserved. Applying
$\tau_\lambda$ at a vertex $v \in V_n$ and using
$W(D_4)$-invariance of $\lambda$
(Definition~\ref{def:fibre-functional}),
\begin{align*}
  S_n^\lambda(g \cdot \Phi)(v)
    &\;=\; \tau_\lambda(\pi^0(g \cdot \Phi))(v)
     \;=\; \lambda\bigl((g \cdot f)(v)\bigr) \\
    &\;=\; \lambda\bigl(\rho_{F^0}^{D_4}(g)\,f(g^{-1}v)\bigr)
     \;=\; \lambda(f(g^{-1}v)) \\
    &\;=\; \tau_\lambda(f)(g^{-1}v)
     \;=\; S_n^\lambda(\Phi)(g^{-1}v)
     \;=\; (g \cdot S_n^\lambda \Phi)(v).
\end{align*}

Statement~(iii) is trivial because $P_n^\PF = \id$ on
$\mathcal T_n = \PF_n$ and the actions on the two identified
spaces are the same.

(iv): Compose (ii), (i), and (iii): $\Pimet^{(n)}(g \cdot \Phi)
= P_n^\PF M_n S_n^\lambda(g \cdot \Phi) = P_n^\PF M_n(g \cdot S_n^\lambda \Phi)
= P_n^\PF(g \cdot M_n S_n^\lambda \Phi) = g \cdot \Pimet^{(n)}(\Phi)$.
\end{proof}

%=====================================================================
\section{Continuity and the Piecewise-Flat Landing Theorem}
\label{sec:continuity}
%=====================================================================

\begin{theorem}[Continuity of $\Pimet^{(n)}$ with explicit bound]
\label{thm:pimet-continuity}
For every $n \geq 0$ the metric projection $\Pimet^{(n)}
\colon X_n \to \PF_n$ is a bounded linear operator, with
\[
  \|\Pimet^{(n)}\|_\op
    \;\leq\; \|P_n^\PF\|_\op \cdot \|M_n\|_\op
       \cdot \|S_n^\lambda\|_\op
    \;\leq\; 1 \cdot |V_n| \cdot \|\lambda\|_\op
    \;=\; |V_n|\,\|\lambda\|_\op.
\]
In particular $\Pimet^{(n)}$ is continuous, and hence Borel
measurable, in the metric topologies on $X_n$ and $\PF_n$.
\end{theorem}

\begin{proof}
$\Pimet^{(n)} = P_n^\PF \circ M_n \circ S_n^\lambda$ is the composition
of three bounded linear operators between finite-dimensional
Hilbert spaces. The operator norm of a composition is at most
the product of the operator norms. Applying the individual
bounds $\|P_n^\PF\|_\op = 1$ (Definition~\ref{def:pn-pf}),
$\|M_n\|_\op \leq |V_n|$ (Theorem~\ref{thm:mn-welldef}(ii)), and
$\|S_n^\lambda\|_\op \leq \|\lambda\|_\op$
(Proposition~\ref{prop:source-bounded}) gives the claimed
estimate. A bounded linear map between metric spaces is
continuous, and every continuous map between metric spaces is
Borel measurable.
\end{proof}

\begin{theorem}[Landing in $\PF_n$]
\label{thm:pimet-landing}
For every $\Phi \in X_n$, $\Pimet^{(n)}(\Phi) \in \PF_n$, i.e.\
$\Pimet^{(n)}(\Phi)$ is a well-defined element of the
piecewise-flat target of Definition~\ref{def:pfn}. The
symmetry admissibility condition --- for every $v \in V_n$,
$\Pimet^{(n)}(\Phi)(v) \in \Sym^2(\R^4)$ --- is satisfied by
construction. No cross-cell compatibility condition is imposed
on $\PF_n$ (Definition~\ref{def:pfn}), so no further check is
required. If the extracted source data $S_n^\lambda(\Phi)$ is
pointwise nonnegative, then $\Pimet^{(n)}(\Phi)(v)$ is positive
semidefinite for every $v$.
\end{theorem}

\begin{proof}
By Theorem~\ref{thm:mn-welldef}(i), $(M_n S_n^\lambda \Phi)(v) \in
\Sym^2(\R^4)$ for every $v$, so
$M_n S_n^\lambda \Phi \in \mathcal T_n = \PF_n$ and the symmetry
admissibility condition holds at every $v$. Applying
$P_n^\PF = \id$ preserves membership in $\mathcal T_n = \PF_n$.
No cross-cell compatibility condition is imposed
(Definition~\ref{def:pfn}). The positive-semidefiniteness
statement under $S_n^\lambda(\Phi) \geq 0$ is
Theorem~\ref{thm:mn-welldef}(iii) composed with $P_n^\PF = \id$.
\end{proof}

%=====================================================================
\section{The Point-Source Rank-One Calculation}
\label{sec:point-source}
%=====================================================================

\begin{definition}[Point source]
\label{def:point-source}
For $w_\ast \in V_n$, the \emph{point source at $w_\ast$} is the
Kronecker function $\delta_{w_\ast} \in \mathcal S_n$ defined by
$\delta_{w_\ast}(w) = 1$ if $w = w_\ast$ and $0$ otherwise.
\end{definition}

\begin{theorem}[Point-source rank-one calculation]
\label{thm:point-source}
Let $w_\ast \in V_n$ and let $\delta_{w_\ast} \in \mathcal S_n$
be the point source at $w_\ast$. Then for every probe vertex
$v \in V_n$ with $v \neq w_\ast$,
\[
  (M_n \delta_{w_\ast})(v)
    \;=\; \hat r_{v, w_\ast}^{(n)}
       \otimes \hat r_{v, w_\ast}^{(n)},
  \qquad
  \hat r_{v, w_\ast}^{(n)}
    \;:=\; \frac{x_{w_\ast}^{(n)} - x_v^{(n)}}
      {\|x_{w_\ast}^{(n)} - x_v^{(n)}\|},
\]
and $(M_n \delta_{w_\ast})(w_\ast) = 0$. For $v \neq w_\ast$, the
$4 \times 4$ symmetric matrix $(M_n \delta_{w_\ast})(v)$ has:
\begin{enumerate}
\item[\textup{(i)}] rank $1$;
\item[\textup{(ii)}] trace $1$;
\item[\textup{(iii)}] spectrum $\{1, 0, 0, 0\}$, i.e.\ eigenvalue
  $1$ with eigenvector $\hat r_{v, w_\ast}^{(n)}$ and eigenvalue
  $0$ with the orthogonal complement as eigenspace;
\item[\textup{(iv)}] traceless part
  $(M_n \delta_{w_\ast})(v) - \tfrac{1}{4}I$ with spectrum
  $\{3/4, -1/4, -1/4, -1/4\}$.
\end{enumerate}
\end{theorem}

\begin{proof}
By Definition~\ref{def:mn},
\[
  (M_n \delta_{w_\ast})(v)
    \;=\; \sum_{w \neq v} \delta_{w_\ast}(w)\,K_n(v, w).
\]
For $v \neq w_\ast$, exactly one term in the sum is nonzero
(namely $w = w_\ast$, which satisfies $w \neq v$ and
$\delta_{w_\ast}(w_\ast) = 1$), giving $(M_n \delta_{w_\ast})(v)
= K_n(v, w_\ast) = \hat r_{v, w_\ast}^{(n)} \otimes
\hat r_{v, w_\ast}^{(n)}$ by Definition~\ref{def:kernel}. For
$v = w_\ast$, the sum ranges over $w \neq w_\ast$ and
$\delta_{w_\ast}(w) = 0$ for every such $w$, so the sum is the
zero matrix.

Claims (i)--(iii) for the case $v \neq w_\ast$ follow from
Lemma~\ref{lem:kernel-properties}(ii) together with the
observation that a rank-one symmetric projector $\hat r \otimes
\hat r$ has $\hat r$ as a unit eigenvector with eigenvalue $1$
(since $(\hat r\,\hat r^\top)\hat r = \hat r\,(\hat r^\top
\hat r) = \hat r$) and the orthogonal hyperplane
$\hat r^\perp$ as a $3$-dimensional eigenspace with eigenvalue
$0$ (since $(\hat r\,\hat r^\top)u = \hat r\,(\hat r^\top u) = 0$
for $u \perp \hat r$).

Claim (iv) follows from (iii) by subtracting $\tfrac{1}{4}I$:
the eigenvalues shift by $-\tfrac{1}{4}$, giving
$\{1 - 1/4, 0 - 1/4, 0 - 1/4, 0 - 1/4\} = \{3/4, -1/4, -1/4,
-1/4\}$.
\end{proof}

\begin{remark}[Finite, not GR]
\label{rmk:not-gr}
Theorem~\ref{thm:point-source} is the theorem-grade finite-level
replacement for the prototype computation in
\texttt{papers/cascade-derivation/cascade-gr.md} \S C4.1. It is a
purely combinatorial identity in symmetric linear algebra on
$\R^4$. No general-relativistic, stress-energy, Einstein-equation,
perfect-fluid, or continuum-metric interpretation is attached to
this theorem, and no such interpretation is endorsed by this
paper.
\end{remark}

%=====================================================================
\section{Self-Audit against the Work Order}
\label{sec:audit}
%=====================================================================

This section checks each Clay-bar success criterion printed in
\S 7 of the work order
\texttt{papers/cascade-metric-projection/WO-CCC-P5.md}.

\begin{itemize}
\item[\textbf{[1]}] \emph{Paper is finite-level only.} Enforced
  in the abstract, \S\ref{sec:intro}, and
  Remark~\ref{rmk:level-bound}. No continuum-limit, convergence,
  GR, or hydrodynamic sentence appears in any theorem or proof.
\item[\textbf{[2]}] \emph{Every symbol is defined before use.}
  Symbols $V_n$, $x_v^{(n)}$ (Notation~\ref{not:vertices-coords});
  $X_n$ (Definition~\ref{def:p4-state}); $\mathcal S_n, S_n^\lambda$
  (Definitions \ref{def:scalar-source} and
  \ref{def:source-extraction}); $K_n(v, w)$
  (Definition~\ref{def:kernel}); $\mathcal T_n$
  (Definition~\ref{def:tn}); $M_n$ (Definition~\ref{def:mn});
  $\PF_n$ (Definition~\ref{def:pfn}); $P_n^\PF$
  (Definition~\ref{def:pn-pf}); $\Pimet^{(n)},
  \widetilde{\Pimet}^{(n)}$ (Definition~\ref{def:pimet});
  $p_{n+1, n}^0, p_{n+1, n}^X$ (Notation~\ref{not:bonding});
  $q_{n+1, n}, q_{n+1, n}^\PF$ (Definition~\ref{def:qn});
  $W(D_4)$ action (Notation~\ref{not:wd4} and
  \ref{not:induced-actions}); $\delta_{w_\ast},
  \hat r_{v, w_\ast}^{(n)}$ (Definition~\ref{def:point-source}
  and Theorem~\ref{thm:point-source}).
\item[\textbf{[3]}] \emph{Moment formula is printed exactly with
  $w \neq v$ exclusion and the squared-norm denominator.}
  Definition~\ref{def:kernel} and Definition~\ref{def:mn}.
\item[\textbf{[4]}] \emph{Denominator nonvanishing is secured under
  an explicit coordinate-injectivity hypothesis on the refinement
  scheme.} Hypothesis~\ref{fact:distinct} (verified at level $0$
  directly and by a relative-interior argument at successor levels,
  with the full Schl\"afli-geometry closure deferred;
  Remark~\ref{rmk:distinct-status}) and
  Corollary~\ref{cor:nonvanishing}. All subsequent results of the
  paper are stated conditional on
  Hypothesis~\ref{fact:distinct}.
\item[\textbf{[5]}] \emph{All of $M_n$, $\PF_n$, $P_n^\PF$,
  $\Pimet^{(n)}$, $q_{n+1, n}$, $W(D_4)$ are printed
  explicitly.} Definitions \ref{def:mn}, \ref{def:pfn},
  \ref{def:pn-pf}, \ref{def:pimet}, \ref{def:qn}, and
  Notations \ref{not:wd4} and \ref{not:induced-actions}.
\item[\textbf{[6]}] \emph{Well-definedness proves boundedness
  and symmetry; it is not a definitional remark.}
  Theorem~\ref{thm:mn-welldef} proves both in explicit
  calculations.
\item[\textbf{[7]}] \emph{Refinement compatibility is an
  operator identity with named maps.}
  Theorem~\ref{thm:mn-refinement} prints
  $q_{n+1, n} \circ M_{n+1} = M_n \circ p_{n+1, n}^0$ with
  both sides named maps defined in \S\ref{sec:refinement}. The
  scope of validity is stated honestly in the theorem statement
  and Remark~\ref{rmk:refinement-scope}: the identity holds on
  level-$(n+1)$ source functions supported on old vertices, and
  is not claimed on the full $\mathcal S_{n+1}$.
\item[\textbf{[8]}] \emph{No sentence upgrades refinement
  stability into convergence, compactness, or continuum
  emergence.} Remark~\ref{rmk:level-bound} and
  Remark~\ref{rmk:refinement-scope} explicitly defer convergence
  to P6.
\item[\textbf{[9]}] \emph{$W(D_4)$-equivariance is proved from
  the P2 coordinate action and the P3 equivariant refinement
  rule, not asserted.} Lemma~\ref{lem:coord-equivariance} proves
  coordinate equivariance $x_{g \cdot v}^{(n)} = g\,x_v^{(n)}$
  for every refined vertex by induction on $n$ from the P2 base
  coordinate action and the P3 midpoint/centroid refinement rule
  (\cite[Def.~2.2]{RefinementSpaces}).
  Lemma~\ref{lem:kernel-equiv} then reduces the equivariance of
  $K_n$ to this coordinate equivariance, and
  Theorem~\ref{thm:mn-equivariance} builds the full theorem from
  the kernel lemma, Definition~\ref{def:fibre-functional} for
  $\lambda$-invariance, and
  \cite[Def.~8.1]{RefinementSpaces} for the state-space action
  (which states unitarity of $\rho_{F^0}^{D_4}(g)$ as part of the
  definition of the branchwise Coxeter action).
\item[\textbf{[10]}] \emph{$\PF_n$ is a concrete
  finite-dimensional Hilbert space with explicit admissibility
  conditions.} Definition~\ref{def:pfn} prints the underlying
  space, the cellwise Frobenius inner product, the induced norm,
  and the admissibility condition (symmetry of each cell tensor).
  No cross-cell compatibility is imposed; this is stated both in
  the definition and in Remark~\ref{rmk:pfn-vs-tn}.
\item[\textbf{[11]}] \emph{Landing theorem checks those
  conditions without hiding behind ``by construction.''}
  Theorem~\ref{thm:pimet-landing} explicitly states that
  symmetry and face-compatibility are met and cites the
  corresponding clauses.
\item[\textbf{[12]}] \emph{Continuity is proved from bounded
  operators with an explicit norm estimate.}
  Theorem~\ref{thm:pimet-continuity} prints the exact
  composition bound.
\item[\textbf{[13]}] \emph{Point-source theorem prints the
  exact rank-one formula and spectrum / trace calculation.}
  Theorem~\ref{thm:point-source} states all four of rank, trace,
  spectrum, and traceless spectrum.
\item[\textbf{[14]}] \emph{Paper never conflates a
  refinement-level metric-side tensor with a continuum metric,
  stress-energy tensor, or Einstein datum.}
  Remark~\ref{rmk:not-gr} states this restriction for
  Theorem~\ref{thm:point-source} specifically, and the abstract
  and \S\ref{sec:intro} state it globally.
\item[\textbf{[15]}] \emph{Citation integrity: no proof cites
  \texttt{cascade-gr.md}, \texttt{foundations.tex}, or scripts.}
  The bibliography below contains only the three cascade
  infrastructure papers P2, P3, P4 and three standard
  mathematics references. No proof cites an informal note.
\item[\textbf{[16]}] \emph{LaTeX hygiene: labels, equation
  labels, notation table, and consistent level superscripts.}
  All theorems and definitions are labelled. The notation audit
  of \S\ref{sec:notation} is consistent with \S\S 2--9. Level
  superscripts $(n)$ and subscripts $n$ are used uniformly.
\end{itemize}

%=====================================================================
\section{Citation Contract and Downstream Handoff}
\label{sec:handoff}
%=====================================================================

\subsection{What downstream papers may cite}

A downstream paper may import from this paper exactly the
following finite-level facts as theorem-grade inputs, with the
stated scope:
\begin{itemize}
\item \textbf{Discrete moment tensor:} the operator
  $M_n \colon \mathcal S_n \to \mathcal T_n$ of
  Definition~\ref{def:mn} and its well-definedness, symmetry,
  and boundedness from Theorem~\ref{thm:mn-welldef}.
\item \textbf{Refinement stability on old-vertex sources:} the
  operator identity $q_{n+1, n} \circ M_{n+1} = M_n \circ
  p_{n+1, n}^0$ of Theorem~\ref{thm:mn-refinement}, valid on
  $\mathcal S_{n+1}^{\mathrm{old}}$.
\item \textbf{$W(D_4)$-equivariance:} the intertwining
  identities $M_n(g\cdot s) = g \cdot M_n s$ and
  $\Pimet^{(n)}(g \cdot \Phi) = g \cdot \Pimet^{(n)}(\Phi)$ of
  Theorem~\ref{thm:mn-equivariance}, for $g \in W(D_4)$.
\item \textbf{Metric projection:} the bounded linear operator
  $\Pimet^{(n)} \colon X_n \to \PF_n$ of
  Definition~\ref{def:pimet} with the explicit norm bound of
  Theorem~\ref{thm:pimet-continuity}.
\item \textbf{Point-source rank-one formula:} the explicit
  identity $M_n \delta_{w_\ast}(v) = \hat r_{v, w_\ast}^{(n)}
  \otimes \hat r_{v, w_\ast}^{(n)}$ of
  Theorem~\ref{thm:point-source}, with the specified rank,
  trace, and spectrum data.
\end{itemize}

\subsection{Recommended citation phrasing}

A downstream paper that wishes to invoke this paper may
write:
\begin{quote}
  ``By Theorem~\ref{thm:mn-welldef} of \cite{MetricProjection},
  the level-$n$ discrete moment tensor $M_n$ is a well-defined
  bounded linear operator from $\ell^2(V_n^{D_4})$ to
  $\ell^2(V_n^{D_4}; \Sym^2(\R^4))$; by
  Theorem~\ref{thm:mn-equivariance} of the same paper, it is
  $W(D_4)$-equivariant; and by Theorem~\ref{thm:mn-refinement}
  it is refinement-compatible on old-vertex source data.''
\end{quote}

\subsection{What this paper does not provide}

We explicitly flag the following as \emph{not} provided by P5,
to avoid downstream over-citation:
\begin{itemize}
\item No convergence of $\{\Pimet^{(n)}\}_{n \geq 0}$ in any
  sense (target of P6).
\item No Gromov--Hausdorff or spectral-stability theorem for the
  refinement family.
\item No continuum metric, Riemannian limit, or Einstein
  equation.
\item No hydrodynamic projection, averaging kernel for a
  macroscopic PDE, or Navier--Stokes statement (target of P7).
\item No full-domain refinement identity on $\mathcal S_{n+1}$
  for all sources. Only the old-vertex restriction is proved.
\item No claim that $P_n^\PF$ is a nontrivial orthogonal
  projection in this paper. It is the identity on
  $\mathcal T_n = \PF_n$ under the admissibility model chosen
  here. Later papers that impose a nontrivial admissibility
  constraint must redefine $P_n^\PF$ and reprove the
  boundedness bound.
\end{itemize}

%=====================================================================
\section{Notation Audit}
\label{sec:notation}
%=====================================================================

\begin{center}
\begin{tabular}{|l|l|l|}
\hline
Symbol & Meaning & First appearance \\
\hline
$V_n$ & level-$n$ $D_4$ refinement vertex set &
  Notation~\ref{not:vertices-coords} \\
$x_v^{(n)}$ & Euclidean coordinate of $v \in V_n$ in $\R^4$ &
  Notation~\ref{not:vertices-coords} \\
$W(D_4)$ & $D_4$ Coxeter group $\leq O(4)$ &
  Notation~\ref{not:wd4} \\
$d_n, D_n$ & minimum / maximum pairwise distance in $V_n$ &
  Notation~\ref{not:diam} \\
$\mathcal S_n$ & scalar source space $\ell^2(V_n)$ &
  Definition~\ref{def:scalar-source} \\
$X_n$ & P4 $D_4$ state space $X_n^{0, D_4} \oplus X_n^{1, D_4}$ &
  Definition~\ref{def:p4-state} \\
$S_n^\lambda$ & source extraction $X_n \to \mathcal S_n$ &
  Definition~\ref{def:source-extraction} \\
$K_n(v, w)$ & normalized rank-one displacement kernel &
  Definition~\ref{def:kernel} \\
$\mathcal T_n$ & raw symmetric-tensor space
  $\ell^2(V_n; \Sym^2(\R^4))$ & Definition~\ref{def:tn} \\
$M_n$ & discrete moment tensor operator
  $\mathcal S_n \to \mathcal T_n$ & Definition~\ref{def:mn} \\
$\PF_n$ & piecewise-flat target
  $\ell^2(V_n; \Sym^2(\R^4))$ & Definition~\ref{def:pfn} \\
$P_n^\PF$ & piecewise-flat projection
  $\mathcal T_n \to \PF_n$ (identity) &
  Definition~\ref{def:pn-pf} \\
$\Pimet^{(n)}$ & metric projection
  $P_n^\PF M_n S_n^\lambda \colon X_n \to \PF_n$ &
  Definition~\ref{def:pimet} \\
$\widetilde{\Pimet}^{(n)}$ & source-level variant
  $P_n^\PF M_n$ & Definition~\ref{def:pimet} \\
$p_{n+1, n}^0$ & P3 vertex bonding (restriction) &
  Notation~\ref{not:bonding} \\
$p_{n+1, n}^X$ & P3 direct-sum state bonding &
  Notation~\ref{not:bonding} \\
$q_{n+1, n}$ & tensor coarse projection (restriction) &
  Definition~\ref{def:qn} \\
$q_{n+1, n}^\PF$ & piecewise-flat coarse projection (restriction)
  & Definition~\ref{def:qn} \\
$\delta_{w_\ast}$ & point source at $w_\ast \in V_n$ &
  Definition~\ref{def:point-source} \\
$\hat r_{v, w_\ast}^{(n)}$ & unit displacement direction &
  Theorem~\ref{thm:point-source} \\
$\mathcal S_{n+1}^{\mathrm{old}}$ & sources supported on
  $V_n \subset V_{n+1}$ &
  Definition~\ref{def:old-vertex-subspace} \\
\hline
\end{tabular}
\end{center}

\begin{thebibliography}{99}

\bibitem{AlgebraicSubstrate}
\emph{Cascade--Classical Correspondence II: The Cascade Algebraic
Substrate},
\texttt{papers/cascade-algebraic-substrate/cascade-algebraic-substrate.tex}.
(Paper P2.) Pinned results used here: \S 3 (the $D_4$
coordinate $4$-plane $U$ and the strict $D_4 \subset E_8$
realization); \S 5 (the $24$-cell coordinate realization of
$V_{24} \subset \R^4$ with Def.~5.3 naming the inscribed
$24$-cell vertex set and Prop.~5.4 together with Cor.~5.5
establishing its regular $24$-cell structure congruent to the
rescaled $D_4$ root polytope).

\bibitem{RefinementSpaces}
\emph{Cascade--Classical Correspondence III: Refinement Spaces
and Symmetry Actions},
\texttt{papers/cascade-refinement-spaces/cascade-refinement-spaces.tex}.
(Paper P3.) Pinned results used here: Def.~2.1 (base $D_4$
graph); Def.~2.2 and Rmk.~2.3 (Schl\"{a}fli refinement with
vertex inclusion $V_n \subset V_{n+1}$); Prop.~2.6 (refinement
properties including finiteness of each $V_n$); Notation~3.3
(sector fibre $F^0$); \S 4 (finite-level $\ell^2$ spaces);
Def.~5.1 (vertex bonding $p^0$ by restriction); Thm.~5.3
(direct-sum bonding and contraction bounds); Def.~8.1 (branchwise
Coxeter action $\rho_{F^0}^\bullet$, including unitarity on each
$X_n^{0, \bullet}$ as part of the definition).

\bibitem{ClosureDynamics}
\emph{Cascade--Classical Correspondence IV: Discrete Closure
Dynamics},
\texttt{papers/cascade-closure-dynamics/cascade-closure-dynamics.tex}.
(Paper P4.) Pinned results used here: Def.~2.1 (finite-level
state space $X_n = X_n^0 \oplus X_n^1$); the direct-sum inner
product on $X_n$; the direct-sum decomposition used in
Definition~\ref{def:source-extraction}.

\bibitem{MetricProjection}
(This paper.) \emph{Cascade--Classical Correspondence V:
$D_4$ Metric Projection from Discrete Moment Tensors}.

\bibitem{ReedSimonI}
M.~Reed and B.~Simon, \emph{Methods of Modern Mathematical
Physics, Vol.~I: Functional Analysis}, rev.\ ed., Academic
Press, 1980. Bounded operators between finite-dimensional
Hilbert spaces, $\ell^2$ theory on finite sets: Ch.~I.1--I.3.

\bibitem{HornJohnson}
R.~A.~Horn and C.~R.~Johnson, \emph{Matrix Analysis}, 2nd ed.,
Cambridge University Press, 2013. Symmetric matrices, rank-one
decomposition, Frobenius inner product, and conjugation by
orthogonal matrices: Ch.~4 and \S 7.7.

\bibitem{Humphreys}
J.~E.~Humphreys, \emph{Reflection Groups and Coxeter Groups},
Cambridge Studies in Advanced Mathematics 29, Cambridge
University Press, 1990. Finite Coxeter groups acting
orthogonally on Euclidean space: Ch.~1.

\end{thebibliography}

\end{document}
